\subsection{Explicit relations and strict examples}
The general recognizer includes two useful identities with short sparse
forms. The following corollary states their hypotheses explicitly.

\begin{corollary}[Two explicit prime-field relations]
\label{sp:thm:constant-field-preprocessing}
Fix a prime $p$. Let a homogeneous sparse list with binary exponents
generate $K\subset S=\Fp[x_0,\ldots,x_N]$, and suppose that
$I=K^{\mathrm{sat}}$ has geometrically integral quotient $S/I$.
For a nonzero listed row $F$, consider the identities
\begin{align}
 F&=U^P+aV^{P-1}U,
       &&a\in\Fp\setminus\{0,-1\},                 \label{sp:eq:constant-two}\\
 F&=U^{P^2}+V^{P^2-P}U^P-V^{P^2-1}U,              \label{sp:eq:constant-three}
\end{align}
where $P=p^r$, $r\ge1$, $U$ is nonzero and homogeneous of degree
$1\le d<P$, and $V$ is any monomial of degree $d$ with coefficient
one. The polynomials $U,V$, the integer $r$, and the scalar $a$ are
not part of the input.

There is a deterministic algorithm of polynomial sparse-bit cost that
finds a valid identity of either displayed type whenever one exists,
and verifies the identity before replacing $F$ by $U$. It may choose
a replacement of minimum degree among all the identities it finds.
Every accepted replacement gives
\[
 K\subseteq K'\subseteq I,\qquad (K')^{\mathrm{sat}}=I.
\]
It introduces no new coordinate index and does not increase the number
of terms. The degree decreases by $P$ in
\eqref{sp:eq:constant-two} and by $P^2$ in
\eqref{sp:eq:constant-three}. Repeated replacements, interleaved with
rank-one digit reduction and any valid coordinate-content divisions,
terminate in polynomial sparse-bit time. The first identity has no
allowed value of $a$ when $p=2$; the second applies in that
characteristic as well.
\end{corollary}
\begin{proof}
For \eqref{sp:eq:constant-two}, use
Theorem~\ref{sp:thm:linearized-preprocessing} with $Q=P$, $s=0$
and $a_0=a$. Its gap condition is $d+1\le P$, and its coefficient
sum is $1+a\ne0$. For \eqref{sp:eq:constant-three}, write $P=p^r$,
take $Q=P^2$, $s=r$, $a_r=1$, $a_0=-1$, and set all other lower
coefficients to zero. The gap again says $d+1\le P$, and the total
coefficient sum is one. This includes characteristic two.
The recognition, minimum-degree choice and iteration bounds are
those of the theorem; neither $V$ nor its membership in the support
of $U$ is required as input.
\end{proof}

\begin{example}[A strict improvement in growing dimension]
\label{sp:ex:constant-field-family}
Fix $p=7$, let $n\ge3$ with $7\nmid n$, and choose $P=7^r>n$.
In $\Fp[x_0,\ldots,x_{n-1}]$ put
\[
 g=\sum_{i=0}^{n-1}x_i^n,\qquad
 F_j=gx_j^{n(P-1)}+g^P\quad(0\le j<n),\qquad K=(F_j).
\]
Then $K^{\mathrm{sat}}=(g)$, and $g$ defines a smooth geometrically
integral anticanonical hypersurface. Every row has $2n-1$ terms and
degree $nP$. No coordinate-content division or nontrivial rank-one
digit reduction changes any row. The least degree of a nonzero
homogeneous equation in $K$ is $nP$, giving least hypersurface excess
$n(P-1)$. Theorem~\ref{sp:thm:linearized-preprocessing} replaces
the rows by scalar multiples of $g$, giving excess zero.
\end{example}
\begin{proof}
The partial derivatives $n x_i^{n-1}$ have no common projective zero.
A reducible projective hypersurface of positive dimension has
intersecting geometric components, so this smooth hypersurface is
geometrically integral. Its degree $n$ is anticanonical in
$\PP^{n-1}$.

Write $F_j=gh_j$, with $h_j=x_j^{n(P-1)}+g^{P-1}$. The identity
\[
 \sum_jx_j^n h_j=2g^P
\]
shows that a common projective zero of the $h_j$ would have $g=0$
and then every $x_j=0$. Thus the cofactors contain a power of the
irrelevant ideal, proving $K^{\mathrm{sat}}=(g)$.
The expansion of $F_j$ has coefficient two at $x_j^{nP}$,
coefficient one at every other pure power, and the $n-1$ distinct
cross terms $x_i^n x_j^{nP-n}$. In particular its monomial content is
one.

Consider a digit radix $Q=7^s$. If $Q\le P$, write $n=aQ+b$,
where $1\le b<Q$. The pure terms have zero low digit and pure high
monomials of degree $nP/Q$. A cross term has low monomial
$x_i^b x_j^{Q-b}$ and high monomial
$x_i^a x_j^{nP/Q-a-1}$, of degree $nP/Q-1$.
The nonzero pure row and a cross row therefore give a nonzero
two-by-two minor with columns of different total degrees. If
$P<Q\le nP$, write $Q=Pq$ and $n=aq+b$, with $a\ge1$ and
$1\le b<q$. Pure terms have distinct low monomials $x_i^{bP}$ and
distinct high monomials $x_i^a$. Cross terms have low monomials
$x_i^n x_j^{bP-n}$ involving two coordinates, since $P>n$.
They do not alter the pure rows, which again have rank at least two.
For $Q>nP$ every high factor is constant. This covers every possible
digit scale.

The ideal is generated in degree $nP$, and a nonzero row attains that
degree. This proves the stated minimum in $K$, not in its saturation.
Identity~\eqref{sp:eq:constant-two} holds with $U=g$, $V=x_j^n$,
and $a=1$. Every accepted replacement belongs to $(g)$, so it has
degree at least $n$. Minimum-degree selection therefore gives a
degree-$n$ associate of $g$. Taking, for example, $n=7r+1$ for
$r\ge2$ gives growing dimension, polynomial written sparse size in
$r$, and exponentially growing former excess.

For any fixed $c\ge2$, adjoin $c-1$ variables $w_i$ and include the
linear rows $w_i$. The saturation becomes $(w_1,\ldots,w_{c-1},g)$
in $\PP^{n+c-2}$, with anticanonical degree sum $n+c-1$.
Before the new reduction, every equation of degree less than $nP$
belongs to the height-$c-1$ linear ideal. A regular $c$-sequence must
therefore have degree sum at least $nP+c-1$, and the linear rows
together with one $F_j$ attain it. After reduction the excess is zero
and the deterministic linear-prefix theorem applies.
\end{proof}

\subsection{Exponential support in a Cartier iterate}
Smooth saturation alone does not bound the number of monomials
reached by repeatedly applying the sparse Cartier operator. The next
family survives coordinate and rank-one digit reduction, while
Theorem~\ref{sp:thm:linearized-preprocessing} recognizes its
prime-field relation. It separates a limitation of explicit support
enumeration from the scalar computation itself.

Multiply-then-section operators also appear in algorithms for
diagonals of rational functions. \citet[Definition~1.7,
Proposition~1.9, Algorithm~1 and Remark~2.2, pp.~249--254]{sp-rowland-yassawi}
construct polynomial numerator states whose finiteness is controlled
by numerical coordinate degrees. The coefficient algorithms of
\citet[Theorems~3.4 and~4.1, author pp.~9--11]{sp-bccd} have
logarithmic dependence on the requested coefficient index, together
with polynomial dependence on numerical degree and height.
\citet[Theorem~1.2 and Corollary~1.4, pp.~4--5; \S2.3]{sp-effective-christol}
bound an invariant linear space by lattice points in a Newton-polytope
region. Those numerical bounds can be large for a sparse list with
binary exponents. Conversely, a monomial-support graph is different
from a minimal automaton or an invariant linear space. The example
below asserts a lower bound only for the specified explicit monomial
representation, and supplies no lower bound for those compressed
representations.

\begin{proposition}[A dense iterate from one monomial]
\label{sp:prop:cartier-states}
Over $\mathbb F_7$, let $P=7^r$, $r\ge1$, and set
\[
 g=x^3+y^3+z^3,\qquad b=3(P-1),\qquad
 f_j=gx_j^b+g^P\quad(x_0=x,x_1=y,x_2=z).
\]
The ideal $K=(f_0,f_1,f_2)$ has smooth geometrically integral
anticanonical saturation $(g)$ and sparse binary input length
$\Theta(r)$. Each cofactor $h_j=f_j/g$ has exactly $28^r$ terms.
Every nonzero constant combination of the rows, even over a field
extension, has monomial content one and has no nontrivial rank-one
base-$7^s$ digit reduction. The least value of $\deg F-3$ for a
nonzero homogeneous $F\in K$ is $b$.

Let $T_f(q)=\mathcal C_7(f_0^6q)$ on homogeneous forms of degree $b$.
Then
\[
 T_f^r((xyz)^{P-1})=(2x^3+y^3+z^3)^{P-1},\qquad
 T_f(h_0)=6h_0.
\]
The first polynomial has exactly $28^r$ nonzero terms. Consequently
an algorithm that explicitly enumerates the reachable monomial states
from $(xyz)^{P-1}$ must handle exponentially many states in the
written sparse input length. This is not a lower bound for computing
the scalar or for using compressed polynomial representations.
\end{proposition}
\begin{proof}
Example~\ref{sp:ex:constant-field-family}, with $n=3$, proves the
saturation, geometry, five-term rows and minimum degree. For nonzero
$c_0,c_1,c_2\in\mathbb F_7$,
\[
 (c_0x^3+c_1y^3+c_2z^3)^{P-1}
  =\prod_{i=0}^{r-1}
    (c_0x^{3\cdot7^i}+c_1y^{3\cdot7^i}
                         +c_2z^{3\cdot7^i})^6.
\]
Each factor has $\binom82=28$ terms, with nonzero multinomial
coefficients modulo seven. After dividing the exponents by three,
their digits at place $i$ are the three chosen nonnegative integers
of sum six. Uniqueness of base-seven expansion prevents collisions
between distinct digit choices. This proves the support count $28^r$.
In $h_j=x_j^b+g^{P-1}$ only one pure coefficient changes from one to
two, so no term disappears.

For a nonzero constant vector $(a_0,a_1,a_2)$, put
$\sigma=a_0+a_1+a_2$. Its row combination is
\[
 F=\sum_j(a_j+\sigma)x_j^{3P}
       +\sum_{j\ne i}a_jx_j^{3P-3}x_i^3.
\]
The pure part is nonzero: its vanishing would give $a_j=-\sigma$
for every $j$ and hence $4\sigma=0$, forcing all coefficients to
vanish. If only one $a_j$ is nonzero, all three pure powers occur.
If at least two are nonzero, then for any coordinate $x_k$ an active
$j\ne k$ and the third index $i$ supply a cross term omitting $x_k$.
Thus the monomial content is one in either case.

At radix $Q=7^s\le P$, write $D=3P/Q$. The zero-residue section is
$\sum_j(a_j+\sigma)X_j^D$. For every active $a_j$, the two
residues $(Q-3)e_j+3e_i$ have sections $a_jX_j^{D-1}$.
These nonzero residues are distinct. If $t$ coefficients $a_j$ are
nonzero, their sections span the $t$ independent monomials
$X_j^{D-1}$, and the nonzero degree-$D$ section is independent of
them. The digit matrix has exact rank $t+1\ge2$.
For $Q>P$, one has $Q\ge7P>3P$, so the high factor is constant.
All these arguments remain valid over extension fields.

The monomial definition of Cartier extraction gives
$\mathcal C_Q(A^QB)=A\mathcal C_Q(B)$ and
$\mathcal C_Q\mathcal C_{Q'}=\mathcal C_{QQ'}$ for powers of seven.
The operator preserves degree $b$, since
$\bigl(6\cdot3P+b-18\bigr)/7=b$.
Induction therefore gives
\[
 T_f^s(q)=\mathcal C_{7^s}(f_0^{7^s-1}q).
\]
Put $A=2x^3+y^3+z^3$ and $B=y^3+z^3$, so that
$f_0=A^P+x^bB$. In the binomial expansion of $f_0^{P-1}$,
a term with $k$ factors from $x^bB$, $1\le k<P$, has $x$-exponent
congruent to $-3k$ modulo $P$, after pulling out its $P$th-power
factor. After multiplication by $(xyz)^{P-1}$ it cannot survive
$\mathcal C_P$, since $P\nmid3k$. The $k=0$ term gives
$A^{P-1}$. This proves the exact iterate, including cancellations.
Its $28^r$ nonzero monomials each have a nonzero transition path in
the matrix product defining $T_f^r$. They are therefore distinct
reachable states. Their ordinary explicit binary exponent lists
require $\Omega(r28^r)$ bits, whereas the input has length
$\Theta(r)$.

Finally,
\[
 T_f(h_0)=\mathcal C_7(g^6h_0^7)
          =h_0\mathcal C_7(g^6).
\]
The only surviving monomial of $g^6$ is $x^6y^6z^6$, with coefficient
$6!/(2!2!2!)=90\equiv6\pmod7$. This proves the eigenvalue assertion.
The equation $g$ and a short expression for the dense iterate are
explicit in this family; their existence is fully compatible with
the lower bound on expanded monomial output.
At this fixed prime and ambient dimension, evaluation of the original
rows at the $57$ points of $\PP^2(\mathbb F_7)$ is itself a
polynomial-time scalar computation. Thus the example concerns the
specified representation of the operator, not the unrestricted scalar
problem.
\end{proof}
